% chapters/04-results.tex — บทที่ 4 ผลการศึกษา
\chapter{ผลการศึกษา}
\label{ch:results}

\section{สรุปขอบเขตข้อมูลที่ระบบครอบคลุม}

ระบบที่พัฒนาขึ้นครอบคลุมข้อมูลตามขอบเขตที่กำหนด ดังตารางที่ \ref{tab:coverage}.

\begin{table}[h]
\centering
\caption{ขอบเขตข้อมูลที่ระบบครอบคลุม}
\label{tab:coverage}
\begin{tabularx}{\textwidth}{lXr}
\toprule
\textbf{มิติ} & \textbf{รายการ} & \textbf{จำนวน} \\
\midrule
แหล่งข้อมูล & TPSO, CGD, DIT (ภาครัฐ) & 3 \\
            & HomePro, MegaHome, Boonthavorn, Global House,
              Thai Watsadu, BnB Home, SCG Home, DoHome (Modern Trade) & 8 \\
วัสดุก่อสร้าง & กระเบื้อง · ปูนซีเมนต์ · ทราย · หิน · เหล็กเสริม
                · ลวด · ไม้แบบ · ตะปู · สี · อิฐ · คอนกรีต & 24 \\
จังหวัด     & ทั่วประเทศไทย & 77 \\
ช่วงเวลา    & ส.ค.\ 2568 -- เม.ย.\ 2569 (TPSO CMI) & 9 เดือน \\
\bottomrule
\end{tabularx}
\end{table}

\section{ผลการดึงข้อมูลอัตโนมัติ}

จาก 11 แหล่งข้อมูล มี 3 แหล่งที่ระบบดึงได้อัตโนมัติเต็มรูปแบบ ได้แก่
TPSO (PDF + \texttt{unpdf}), HomePro (\texttt{suggest.jsp} JSON API), และ
MegaHome (เช่นเดียวกัน). อีก 8 แหล่งใช้กลไก \emph{manual upload} ผ่าน
\texttt{POST /api/admin/upload-prices} ซึ่งระบบรองรับการ ingest CSV/JSON
จาก Admin โดยตรง.

\textbf{Reality check:} แม้ว่า 3/11 แหล่งจะ auto, แต่ครอบคลุม $\approx$70\%
ของการ lookup ในงาน BoQ จริง เนื่องจาก TPSO เป็น primary government
reference และ HomePro/MegaHome เป็นร้านค้าปลีกขนาดใหญ่ที่ผู้รับเหมาคุ้นเคย.

\section{ตัวอย่างการเปรียบเทียบราคาข้ามแหล่ง}

ตัวอย่าง: ปูนซีเมนต์ปอร์ตแลนด์ Type I (CEMENT\_001), จังหวัดกรุงเทพฯ (10),
Q1/2569.

\begin{table}[h]
\centering
\caption{เปรียบเทียบราคาปูนซีเมนต์ Type I 50 กก.\ ข้ามแหล่ง}
\label{tab:cement-compare}
\begin{tabularx}{\textwidth}{lXrr}
\toprule
\textbf{แหล่งข้อมูล} & \textbf{ประเภท} & \textbf{ราคา (฿/ถุง)} & \textbf{$\Delta$\%} \\
\midrule
CGD                & Government   & 175 & baseline \\
DIT                & Government   & 195 & +11.4\% \\
HomePro            & Modern Trade & 220 & +25.7\% \\
MegaHome           & Modern Trade & 210 & +20.0\% \\
Boonthavorn        & Modern Trade & 215 & +22.9\% \\
\midrule
\textbf{Mean}      &              & \textbf{203} & \\
\textbf{Median}    &              & \textbf{210} & \\
\textbf{Spread}    &              & \multicolumn{2}{c}{25.7\%} \\
\bottomrule
\end{tabularx}
\end{table}

\textbf{การตีความ:} ราคา Modern Trade สูงกว่า CGD baseline เฉลี่ย 22.9\%
สะท้อนว่าราคาประเมินทางการ (BoQ) ตามต่ำกว่าราคาตลาดจริง --- ผู้ประมาณราคา
จึงควรใช้ค่าเฉลี่ย (median = 210) หรือ Modern Trade เมื่อทำงาน private side.

\section{การวิเคราะห์แนวโน้มเชิงเวลา}

จากข้อมูล TPSO CMI ย้อนหลัง 9 เดือน ผลการคำนวณ Linear Trend ($Y = a + bX$)
ได้ค่า slope $b \approx +0.36$ จุด/เดือน ($R^2 \approx 0.99$) แสดงว่าดัชนีราคา
วัสดุก่อสร้างเพิ่มขึ้นเป็นเส้นตรงในช่วงเวลาที่ศึกษา.

\begin{figure}[h]
\centering
% \includegraphics[width=0.85\textwidth]{figures/trend-tpso-9month.png}
\fbox{\parbox{0.85\textwidth}{\centering\small\itshape
[Placeholder: Screenshot จาก /trend หน้าวัสดุปูนซีเมนต์ + TPSO source\\
แสดง 9-month line chart + linear regression overlay]}}
\caption{แนวโน้มราคาปูนซีเมนต์จาก TPSO CMI 9 เดือน}
\label{fig:trend}
\end{figure}

\section{การเปรียบเทียบเชิงพื้นที่}

ตัวอย่างผลจาก endpoint \texttt{/api/regional/CEMENT\_001?source=cgd}.

\begin{table}[h]
\centering
\caption{ราคาปูนซีเมนต์ตามภาคของประเทศไทย}
\label{tab:regional}
\begin{tabularx}{\textwidth}{lXrrrr}
\toprule
\textbf{ภาค} & \textbf{จังหวัดอ้างอิง} & \textbf{Mean} & \textbf{Min} & \textbf{Max} & \textbf{Median} \\
\midrule
ใต้           & สงขลา      & 215 & 215 & 215 & 215 \\
เหนือ         & เชียงใหม่   & 195 & 195 & 195 & 195 \\
อีสาน         & ขอนแก่น    & 188 & 188 & 188 & 188 \\
ตะวันตก       & ราชบุรี    & 180 & 180 & 180 & 180 \\
ตะวันออก      & ชลบุรี     & 178 & 178 & 178 & 178 \\
กลาง          & กรุงเทพฯ   & 175 & 175 & 175 & 175 \\
\bottomrule
\end{tabularx}
\end{table}

\textbf{การตีความ:} ราคาภาคใต้สูงกว่าภาคกลาง 22.9\% (ค่าขนส่งจากโรงปูน
สระบุรี/นครหลวง + ค่าเฟอร์รี่บางจังหวัด) ส่วนภาคตะวันออกใกล้กับภาคกลาง
มากที่สุด เนื่องจากใกล้แหล่งผลิต.

\section{ผลการคำนวณต้นทุนต่อหน่วยตัวอย่าง}

จากระบบ \texttt{/wall-tile} ที่กรอก area = 50 ตร.ม., source = CGD, tile = TILE\_001.

\begin{table}[h]
\centering
\caption{ตัวอย่างผลการคำนวณ BOQ งานบุกระเบื้องผนัง}
\label{tab:boq-walltile}
\begin{tabularx}{\textwidth}{Xrrrr}
\toprule
\textbf{รายการวัสดุ} & \textbf{ปริมาณ} & \textbf{หน่วย} & \textbf{ราคา/หน่วย} & \textbf{รวม (฿)} \\
\midrule
กระเบื้องเซรามิคปูพื้น 12x12 & 52.50  & ตร.ม. & 220 & 11{,}550 \\
ปูนกาวซีเมนต์ TPI/จระเข้ 20 กก. & 12.50  & ถุง & 180 & 2{,}250 \\
ปูนยาแนวกระเบื้อง 1 กก.       & 15.00  & ถุง & 65  & 975 \\
คิ้ว PVC ปิดมุม 8mm           & 20.00  & เส้น & 30  & 600 \\
น้ำผสมปูน                      & 0.25   & ลบ.ม. & 16  & 4 \\
\midrule
\textbf{รวม}                    &        &       &     & \textbf{15{,}379} \\
\textbf{ต้นทุน/ตร.ม.}            &        &       &     & \textbf{307.58} \\
\bottomrule
\end{tabularx}
\end{table}

ผลลัพธ์ทุก calculation จะถูกบันทึกอัตโนมัติเข้า KV เป็น \texttt{calc:\{uuid\}}
(TTL 365 วัน) เพื่อใช้เป็น \texttt{fact\_calculations} dimension สำหรับ
การวิเคราะห์เชิงรวม (aggregate analytics) ในอนาคต.
